Las pendientes se anulan sobre \(y=t^2\), son positivas debajo y negativas encima. La solución que parte de cero queda debajo para \(t>0\) y crece. La que parte de dos baja hasta \((1,1)\), cruza allí la parábola con tangente horizontal y después crece. Las soluciones no se cruzan entre sí. La parábola de pendientes nulas no es una solución ni una barrera.
